\documentclass[12pt]{article}
\usepackage[utf8]{inputenc}
\usepackage{amsmath}
\usepackage{amssymb}
\usepackage{graphicx}
\usepackage{booktabs}
\usepackage{listings}
\usepackage{xcolor}
\usepackage{hyperref}
\usepackage{geometry}
\geometry{margin=1in}

\lstset{
    language=Python,
    basicstyle=\ttfamily\small,
    keywordstyle=\color{blue}\bfseries,
    commentstyle=\color{green!60!black},
    stringstyle=\color{red},
    numbers=left,
    numberstyle=\tiny\color{gray},
    stepnumber=1,
    numbersep=5pt,
    backgroundcolor=\color{gray!10},
    frame=single,
    breaklines=true,
    breakatwhitespace=false,
    tabsize=4,
    showstringspaces=false
}

\title{Utility Modules Documentation}
\author{}
\date{}

\begin{document}

\maketitle

\section{Introduction}

This document describes the utility modules used throughout the CTCC: \texttt{yaml\_loaders.py}, \texttt{financial\_utils.py}, \texttt{calculation\_utils.py}, and \texttt{weighted\_miles.py}. These modules provide shared functionality for loading configuration data, performing financial calculations, computing transmission line parameters, and calculating terrain-adjusted distances.

\section{yaml\_loaders.py}

\subsection{Method Overview}

The \texttt{yaml\_loaders.py} module provides centralized functions for loading YAML configuration files. It eliminates code duplication by providing a single interface for accessing project parameters, financing details, physical details, and cost parameters.

\subsection{Key Functions}

\subsubsection{\texttt{load\_project\_technical\_details()}}

Loads project technical details from \texttt{01\_project\_technical\_details.yaml}.

\textbf{Returns:}
\begin{itemize}
    \item Tuple: (construction\_type, ac\_dc, capacity\_mw, conductor\_type, converter\_type, line\_utilization, reconductoring, delay\_years, construction\_years, project\_lifetime)
\end{itemize}

\subsubsection{\texttt{load\_physical\_details()}}

Loads physical project details and returns total miles only.

\textbf{Returns:}
\begin{itemize}
    \item \texttt{total\_miles} (float): Sum of all terrain miles
\end{itemize}

\subsubsection{\texttt{load\_physical\_details\_detailed()}}

Loads physical project details with detailed terrain breakdown.

\textbf{Returns:}
\begin{itemize}
    \item Tuple: (total\_miles, forested\_miles, scrubbed\_flat\_miles, wetland\_miles, farmland\_miles, desert\_barren\_miles, urban\_miles, rolling\_hills\_miles, mountain\_miles, subsea\_miles)
\end{itemize}

\subsubsection{\texttt{load\_financing\_details()}}

Loads financing parameters and calculates real WACC using Fisher equation.

\textbf{Returns:}
\begin{itemize}
    \item Tuple: (inflation\_rate, base\_year, wacc\_nominal, wacc\_real)
\end{itemize}

\textbf{Key Logic:}
\begin{lstlisting}
wacc_real = (1 + wacc_nominal) / (1 + inflation_rate) - 1
\end{lstlisting}

\subsubsection{\texttt{load\_circuit\_and\_resistance\_details(category)}}

Loads circuit and resistance details for specified project category.

\textbf{Parameters:}
\begin{itemize}
    \item \texttt{category} (str): Project category identifier
\end{itemize}

\textbf{Returns:}
\begin{itemize}
    \item Tuple: (voltage\_kv, conductors\_per\_phase, number\_of\_phases, number\_of\_circuits\_poles, AC\_75\_resistance, DC\_20\_resistance)
\end{itemize}

\subsubsection{Other Loader Functions}

Additional loader functions include:
\begin{itemize}
    \item \texttt{load\_row\_widths(category)} - ROW width for project category
    \item \texttt{load\_row\_details()} - ROW cost parameters
    \item \texttt{load\_delay\_costs()} - Delay cost components
    \item \texttt{load\_emissions\_details()} - Emissions parameters
    \item \texttt{load\_congestion\_reductions()} - Congestion reduction parameters
    \item \texttt{load\_curtailment\_reductions()} - Curtailment reduction parameters
    \item \texttt{load\_contingencies()} - Contingency percentages
    \item \texttt{load\_environmental\_mitigation()} - Environmental mitigation parameters
    \item \texttt{load\_cost\_timing\_patterns()} - AFUDC timing patterns
    \item \texttt{load\_afudc\_config()} - AFUDC configuration
    \item \texttt{load\_insurance\_details()} - Insurance parameters
    \item \texttt{load\_wildfire\_costs()} - Wildfire parameters
    \item \texttt{load\_outage\_costs()} - Outage parameters
\end{itemize}

\subsection{Example}

\begin{lstlisting}
from yaml_loaders import (
    load_project_technical_details,
    load_financing_details,
    load_physical_details
)

# Load project specs
construction_type, ac_dc, capacity_mw, ... = load_project_technical_details()

# Load financing
inflation_rate, base_year, wacc_nominal, wacc_real = load_financing_details()

# Load physical details
total_miles = load_physical_details()
\end{lstlisting}

\section{financial\_utils.py}

\subsection{Method Overview}

The \texttt{financial\_utils.py} module provides shared financial calculation functions for present value calculations, amortization, and AFUDC capitalization.

\subsection{Key Functions}

\subsubsection{\texttt{calculate\_present\_value(annual\_cost, wacc\_real, total\_years, start\_year=1)}}

Calculates the present value of annual payments over a given time period.

\textbf{Parameters:}
\begin{itemize}
    \item \texttt{annual\_cost} (float): Annual cost amount
    \item \texttt{wacc\_real} (float): Real WACC (discount rate)
    \item \texttt{total\_years} (float): Number of years
    \item \texttt{start\_year} (int): Year when payments begin (default: 1)
\end{itemize}

\textbf{Returns:}
\begin{itemize}
    \item \texttt{present\_value} (float): Present value of the payment stream
\end{itemize}

\textbf{Formula:}
\[
PV = \sum_{t=start}^{start+total-1} \frac{C}{(1 + r)^t}
\]

\subsubsection{\texttt{calculate\_amortized\_cost(principal, wacc\_real, project\_lifetime)}}

Calculates annual amortized cost using standard amortization formula.

\textbf{Parameters:}
\begin{itemize}
    \item \texttt{principal} (float): Initial cost
    \item \texttt{wacc\_real} (float): Real WACC (discount rate)
    \item \texttt{project\_lifetime} (int): Number of years to amortize
\end{itemize}

\textbf{Returns:}
\begin{itemize}
    \item \texttt{annual\_payment} (float): Annual amortized payment
\end{itemize}

\textbf{Formula:}
\[
A = P \times \frac{r(1+r)^n}{(1+r)^n - 1}
\]

where $P$ is principal, $r$ is rate, and $n$ is number of years.

\subsubsection{\texttt{calculate\_afudc\_rate(financing\_yaml)}}

Calculates AFUDC rate from capital structure or falls back to WACC.

\textbf{Parameters:}
\begin{itemize}
    \item \texttt{financing\_yaml} (dict): Loaded financing YAML data
\end{itemize}

\textbf{Returns:}
\begin{itemize}
    \item Tuple: (\texttt{afudc\_rate}, \texttt{source\_description})
\end{itemize}

\textbf{Key Logic:}
\begin{lstlisting}
if capital_structure specified:
    rate = equity_percent * cost_of_equity + debt_percent * cost_of_debt
    source = "capital structure"
else:
    rate = wacc_nominal
    source = "WACC nominal"
\end{lstlisting}

\subsubsection{\texttt{calculate\_afudc\_capitalized\_cost(nominal\_cost, timing\_pattern, delay\_years, construction\_years, afudc\_rate, delay\_has\_active\_work)}}

Capitalizes a cost using AFUDC (compounds forward to Commercial Operation Date).

\textbf{Parameters:}
\begin{itemize}
    \item \texttt{nominal\_cost} (float): Total nominal cost
    \item \texttt{timing\_pattern} (dict): Dict with \texttt{during\_delay}, \texttt{during\_construction}, \texttt{afudc\_eligible}
    \item \texttt{delay\_years} (float): Number of delay years
    \item \texttt{construction\_years} (float): Number of construction years
    \item \texttt{afudc\_rate} (float): AFUDC rate (annual)
    \item \texttt{delay\_has\_active\_work} (bool): Whether active CWIP work continues during delay
\end{itemize}

\textbf{Returns:}
\begin{itemize}
    \item Tuple: (\texttt{capitalized\_cost}, \texttt{afudc\_amount})
\end{itemize}

\textbf{Key Logic:}
\begin{lstlisting}
# Cost during delay period
delay_cost = nominal_cost * timing_pattern["during_delay"]
if delay_has_active_work:
    # Compound from midpoint of delay to COD
    avg_years_to_cod = total_years_to_cod - (delay_years / 2)
    capitalized_cost += delay_cost * (1 + afudc_rate) ** avg_years_to_cod
else:
    # No AFUDC during delay, compound only during construction
    capitalized_cost += delay_cost * (1 + afudc_rate) ** construction_years

# Cost during construction period
construction_cost = nominal_cost * timing_pattern["during_construction"]
# Compound from midpoint of construction to COD
avg_years_to_cod = construction_years / 2
capitalized_cost += construction_cost * (1 + afudc_rate) ** avg_years_to_cod

afudc_amount = capitalized_cost - nominal_cost
\end{lstlisting}

\subsection{Example}

\begin{lstlisting}
from financial_utils import (
    calculate_present_value,
    calculate_amortized_cost,
    calculate_afudc_capitalized_cost
)

# Present value of annual O&M costs
pv = calculate_present_value(
    annual_cost=1000000,
    wacc_real=0.03,
    total_years=50,
    start_year=3
)

# Amortize build cost
annual_payment = calculate_amortized_cost(
    principal=100000000,
    wacc_real=0.03,
    project_lifetime=50
)

# AFUDC capitalization
capitalized, afudc = calculate_afudc_capitalized_cost(
    nominal_cost=100000000,
    timing_pattern={"during_delay": 0.1, "during_construction": 0.9, "afudc_eligible": True},
    delay_years=2,
    construction_years=2,
    afudc_rate=0.05,
    delay_has_active_work=False
)
\end{lstlisting}

\section{calculation\_utils.py}

\subsection{Method Overview}

The \texttt{calculation\_utils.py} module provides shared calculation functions for transmission line parameters, including phase current, full load adjustment, line losses, and converter losses.

\subsection{Key Functions}

\subsubsection{\texttt{calculate\_phase\_current(capacity\_mw, voltage\_kv, number\_of\_phases, number\_of\_circuits\_poles, ac\_dc)}}

Calculates phase current based on AC or DC transmission type.

\textbf{Parameters:}
\begin{itemize}
    \item \texttt{capacity\_mw} (float or str): Line capacity in MW
    \item \texttt{voltage\_kv} (float): Voltage in kV
    \item \texttt{number\_of\_phases} (int): Number of phases
    \item \texttt{number\_of\_circuits\_poles} (int): Number of circuits/poles
    \item \texttt{ac\_dc} (str): "AC" or "DC"
\end{itemize}

\textbf{Returns:}
\begin{itemize}
    \item \texttt{phase\_current} (float): Phase current in Amps
\end{itemize}

\textbf{Formulas:}
\begin{align}
I_{AC} &= \frac{P \times 1000}{0.95 \times V \times \sqrt{N_{phases}} \times N_{circuits}} \\
I_{DC} &= \frac{P \times 1000}{V \times \sqrt{N_{phases}} \times N_{poles}}
\end{align}

\subsubsection{\texttt{full\_load\_adjusted(line\_utilization\_percent)}}

Calculates full load adjustment based on line utilization.

\textbf{Parameters:}
\begin{itemize}
    \item \texttt{line\_utilization\_percent} (float): Line utilization as decimal (0-1)
\end{itemize}

\textbf{Returns:}
\begin{itemize}
    \item \texttt{adjustment\_factor} (float): Full load adjustment factor
\end{itemize}

\textbf{Formula:}
\[
\text{Adjustment} = \frac{u + u^2}{2}
\]

where $u$ is the utilization factor.

\subsubsection{\texttt{calculate\_line\_losses(...)}}

Calculates line losses for transmission lines.

\textbf{Parameters:}
\begin{itemize}
    \item \texttt{phase\_current} (float): Phase current in Amps
    \item \texttt{full\_load\_adj} (float): Full load adjustment factor
    \item \texttt{AC\_75\_resistance} (float): AC resistance at 75°C (ohms/mile)
    \item \texttt{DC\_20\_resistance} (float): DC resistance at 20°C (ohms/mile)
    \item \texttt{ac\_dc} (str): "AC" or "DC"
    \item \texttt{number\_of\_circuits\_poles} (int): Number of circuits/poles
    \item \texttt{conductors\_per\_phase} (int): Number of conductors per phase
    \item \texttt{number\_of\_phases} (int): Number of phases
    \item \texttt{line\_length} (float): Line length in miles
    \item \texttt{project\_lifetime} (int): Project lifetime in years
    \item \texttt{capacity\_mw\_numeric} (float): Capacity in MW
    \item \texttt{line\_utilization\_percent} (float): Line utilization as decimal
\end{itemize}

\textbf{Returns:}
\begin{itemize}
    \item Tuple: (losses\_mwh\_per\_year, lifetime\_losses\_mwh, losses\_mw\_per\_mile, resistance\_per\_mile, line\_loss\_per\_mile\_percent, total\_line\_loss\_mw, total\_line\_loss\_percent)
\end{itemize}

\textbf{Key Logic:}
\begin{lstlisting}
# Select resistance based on AC/DC
resistance_per_mile = AC_75_resistance if ac_dc == "AC" else DC_20_resistance

# Calculate losses per mile
number_of_conductors = conductors_per_phase * number_of_phases * number_of_circuits_poles
losses_mw_per_mile = ((phase_current / conductors_per_phase) ** 2 
                     * number_of_conductors 
                     * resistance_per_mile 
                     * full_load_adj / 1000000)

# Total losses
total_line_loss_mw = losses_mw_per_mile * line_length
losses_mwh_per_year = total_line_loss_mw * 8760
lifetime_losses_mwh = losses_mwh_per_year * project_lifetime
\end{lstlisting}

\subsubsection{\texttt{calculate\_converter\_losses(number\_of\_converters, converter\_type, line\_utilization\_percent, capacity\_mw\_numeric, ac\_dc)}}

Calculates converter losses based on number of converters (DC only).

\textbf{Parameters:}
\begin{itemize}
    \item \texttt{number\_of\_converters} (int): Number of converters
    \item \texttt{converter\_type} (str): Converter type (LCC or VSC)
    \item \texttt{line\_utilization\_percent} (float): Line utilization as decimal
    \item \texttt{capacity\_mw\_numeric} (float): Capacity in MW
    \item \texttt{ac\_dc} (str): "AC" or "DC"
\end{itemize}

\textbf{Returns:}
\begin{itemize}
    \item Tuple: (total\_converter\_losses\_mw, total\_converter\_losses\_mwh, converter\_loss\_percent)
\end{itemize}

\textbf{Key Logic:}
\begin{lstlisting}
if ac_dc == "DC":
    converter_loss_percentage = 0.0075 if "LCC" in converter_type else 0.01
    converter_losses_mw = (converter_loss_percentage 
                          * line_utilization_percent 
                          * capacity_mw_numeric)
    total_converter_losses_mw = converter_losses_mw * number_of_converters
    total_converter_losses_mwh = total_converter_losses_mw * 8760
else:
    total_converter_losses_mw = 0
    total_converter_losses_mwh = 0
\end{lstlisting}

\subsection{Example}

\begin{lstlisting}
from calculation_utils import (
    calculate_phase_current,
    full_load_adjusted,
    calculate_line_losses
)

# Calculate phase current
phase_current = calculate_phase_current(
    capacity_mw=500,
    voltage_kv=345,
    number_of_phases=3,
    number_of_circuits_poles=1,
    ac_dc="AC"
)

# Calculate full load adjustment
full_load_adj = full_load_adjusted(0.7)  # 70% utilization

# Calculate line losses
losses_mwh_per_year, lifetime_losses_mwh, ... = calculate_line_losses(
    phase_current, full_load_adj, AC_75_resistance, DC_20_resistance,
    "AC", 1, 2, 3, 100, 50, 500, 0.7
)
\end{lstlisting}

\section{weighted\_miles.py}

\subsection{Method Overview}

The \texttt{weighted\_miles.py} module calculates weighted miles and average terrain multiplier for transmission lines. Weighted miles account for terrain complexity by multiplying miles in each terrain type by a terrain-specific cost multiplier.

\subsection{Key Functions}

\subsubsection{\texttt{calculate\_weighted\_miles()}}

Calculates weighted miles and average terrain multiplier from project physical details.

\textbf{Returns:}
\begin{itemize}
    \item Tuple: (\texttt{weighted\_miles}, \texttt{average\_terrain\_multiplier})
\end{itemize}

\textbf{Calculation:}
\begin{align}
\text{Weighted Miles} &= \sum_{\text{terrains}} \text{Miles}_{terrain} \times \text{Multiplier}_{terrain} \\
\text{Average Multiplier} &= \frac{\text{Weighted Miles}}{\text{Total Miles}}
\end{align}

\textbf{Key Logic:}
\begin{lstlisting}
weighted_miles = 0
for terrain_type, terrain_miles in terrain_miles_dict.items():
    multiplier = terrain_multipliers[terrain_type]
    weighted_miles += terrain_miles * multiplier

average_terrain_multiplier = weighted_miles / total_miles
\end{lstlisting}

\subsection{Example}

\begin{lstlisting}
from weighted_miles import calculate_weighted_miles

# Calculate weighted miles
weighted_miles, avg_multiplier = calculate_weighted_miles()

# Example calculation:
# Terrain distribution:
#   - Forested: 40 miles, multiplier: 2.25
#   - Farmland: 60 miles, multiplier: 1.0
#   - Total: 100 miles
#
# Weighted miles = (40 * 2.25) + (60 * 1.0) = 150 miles
# Average multiplier = 150 / 100 = 1.5
\end{lstlisting}

\section{Integration Examples}

\subsection{Combined Usage}

\begin{lstlisting}
# Typical usage in a cost calculation script:
from yaml_loaders import (
    load_project_technical_details,
    load_financing_details,
    load_physical_details
)
from financial_utils import calculate_present_value
from weighted_miles import calculate_weighted_miles

# Load configuration
construction_type, ac_dc, capacity_mw, ... = load_project_technical_details()
inflation_rate, base_year, wacc_nominal, wacc_real = load_financing_details()
total_miles = load_physical_details()

# Calculate weighted miles
weighted_miles, avg_multiplier = calculate_weighted_miles()

# Calculate present value
pv = calculate_present_value(
    annual_cost=1000000,
    wacc_real=wacc_real,
    total_years=50,
    start_year=3
)
\end{lstlisting}

\section{Key Implementation Details}

\subsection{Path Resolution}

All YAML loaders use relative paths (\texttt{"../yamls/"}) assuming scripts are run from the project root directory. The \texttt{csv\_output\_manager.py} uses dynamic path resolution to handle different execution contexts.

\subsection{Fisher Equation}

The real WACC calculation uses the Fisher equation to convert nominal WACC to real WACC, accounting for inflation:
\[
r_{real} = \frac{1 + r_{nominal}}{1 + \pi} - 1
\]

\subsection{Full Load Adjustment}

The full load adjustment accounts for the fact that losses are proportional to current squared, and average current is less than peak current. The formula $(u + u^2)/2$ provides a better approximation than simple $u^2$ for typical utilization patterns.

\subsection{AFUDC Timing}

AFUDC capitalization compounds costs forward to Commercial Operation Date (COD). The calculation accounts for:
\begin{itemize}
    \item Costs during delay period (may or may not have AFUDC depending on active work)
    \item Costs during construction period (always has AFUDC)
    \item Uniform spending assumption (uses midpoint of each period)
\end{itemize}

\subsection{Terrain Multipliers}

Terrain multipliers reflect the relative cost/complexity of construction in different terrain types. For example, forested terrain may have a multiplier of 2.25, meaning construction costs 2.25x more per mile than baseline terrain.

\end{document}

